\section{Trabalhos Relacionados}
\label{sec:trabalhos_relacionados}

Esta seção apresenta estudos empíricos relacionados à segurança, testes, qualidade e geração automatizada de artefatos de Infraestrutura como Código, destacando o método, os principais resultados e a relação de cada um com a proposta desta pesquisa.

Rahman et al. \cite{rahman2021securitysmells} investigaram \textit{security smells} em scripts Ansible e Chef com a ferramenta SLAC, identificando 46.600 ocorrências, incluindo 7.849 senhas codificadas diretamente nos scripts, evidência de que o uso de \textit{IaC} não elimina problemas de segurança. Verdet et al. \cite{verdet2025security} analisaram 812 projetos Terraform \textit{open source} com Checkov e Tfsec, observando maior adoção de políticas de controle de acesso do que de criptografia em repouso, o que reforça que a existência de recomendações de segurança não garante sua aplicação. Diferentemente desses trabalhos, esta pesquisa não analisa apenas artefatos existentes, mas compara \textit{templates} produzidos manualmente e por IA em duas condições de contexto, usando análise estática em arquivos Terraform como um dos instrumentos de avaliação.

Hassan et al. \cite{hassan2024state} estudaram defeitos de reconciliação de estado em \textit{IaC} (divergência entre o estado real e o esperado), identificando oito categorias de falhas com apoio de modelos de linguagem. Os resultados mostram que a ausência de erros sintáticos não garante que o estado produzido corresponda ao esperado, o que sustenta a inclusão, neste trabalho, de verificações de execução além das análises estáticas. Sokolowski, Spielmann e Salvaneschi \cite{sokolowski2024automated} propuseram a \textit{Automated Configuration Testing} para programas de \textit{IaC} de propósito geral, avaliando 6.081 programas entre projetos do GitHub e \textit{benchmarks} artificiais. Enquanto os autores propõem uma técnica específica de teste, esta pesquisa utiliza pipelines de \textit{CI/CD} para aplicar um conjunto comum de verificações a artefatos produzidos por diferentes métodos.

Zhang et al. \cite{zhang2024generative} avaliaram 98 manifestos Kubernetes gerados pelo ChatGPT e identificaram pelo menos um \textit{smell} em 35,8\% deles, sobretudo ausência de limites de CPU e memória, concluindo que modelos generativos produzem configurações funcionais que ainda exigem garantia de qualidade. Esta pesquisa amplia essa perspectiva ao utilizar Terraform e comparar geração manual, IA sem contexto organizacional e IA com contexto organizacional.

Em conjunto, os estudos mostram que artefatos de infraestrutura estão sujeitos a vulnerabilidades, defeitos de estado e problemas de qualidade, e que mecanismos automatizados são relevantes para sua avaliação, cada qual com um foco distinto, seja segurança, defeitos de estado, técnicas de teste ou geração por IA. Nenhum deles, porém, combina em um mesmo experimento desenvolvimento manual, geração por IA sem contexto e geração por IA com contexto organizacional, avaliando simultaneamente eficiência operacional, conformidade, padronização e segurança, lacuna que esta pesquisa busca preencher.
